\kbvbooktitle

\kbvchapter{ಪಾಠ ೧}{ವಿಷಯ ಪ್ರವೇಶ}

\indent ನಮ್ಮ ಮನಸ್ಸಿನಲ್ಲಿರುವ ಅಭಿಪ್ರಾಯವನ್ನು ಮತ್ತೊಬ್ಬರಿಗೆ ತಿಳಿಸಲು ನೆರವಾಗುವ ಒಂದು ಸಾಧನವೇ ಭಾಷೆ. ನಾವು ಆಡು ಭಾಷೆಯಲ್ಲಿ ಅಭಿಪ್ರಾಯವ್ಯಕ್ತಪಡಿಸುವ ಶಬ್ದರೂಪ ಮಾಧ್ಯಮವೇ ಮಾತು. ಉದಾಹರಣೆಗೆ, “ಕಮಲಳ ಮಗ ಬಸವ”-ಈ ಮಾತಿನ ಅರ್ಥವನ್ನು ನಾವು ಸುಲಭವಾಗಿ ಗ್ರಹಿಸಬಲ್ಲೆವು. ಕಮಲ ಯಾರು? ಬಸವ ಯಾರು? ಅವರಿಬ್ಬರಿಗಿರುವ ಸಂಬಂಧವೇನು? ಎಂಬುದು ನಮಗೆ ತಿಳಿದು ಬರುವುದು; ಇದು ಅರ್ಥಪೂರ್ಣವಾದ ಮಾತು. ಹೀಗೆ ಪೂರ್ಣವಾದ ಅರ್ಥವನ್ನು ಕೊಡುವ ಮಾತಿಗೆ ವಾಕ್ಯ ಎನ್ನುತ್ತೇವೆ.\par

\indent ಮೇಲೆ ಉದಾಹರಿಸಿದ ವಾಕ್ಯದಲ್ಲಿ ಮೂರು ಭಾಗಗಳಿವೆ ; ಕಮಲಳ ಮಗ ಬಸವ-ಇವೇ ಆ ಮೂರು ಭಾಗಗಳು. ಇವುಗಳಲ್ಲಿ ಯಾವುದೆ ಒಂದನ್ನು ಹೇಳಿದರೂ ಮಾತಿನ ಅರ್ಥ ಬರುವುದಿಲ್ಲ; ಇಂತಹ ಬಿಡಿ ಭಾಗಗಳಿಗೆ “ಪದ”ಗಳೆಂದು ಹೆಸರು.\par

\indent ಮೇಲಿನ ಮೂರು ಪದಗಳಲ್ಲಿ ಒಂದನ್ನು ಗಮನಿಸೋಣ ; ಉದಾಹರಣೆಗೆ, “ಬಸವನು” ಈ ಪದದಲ್ಲಿ ಬಿಡಿಯಾದ ಮೂರು ಅಕ್ಷರಗಳಿವೆ. ಅಕ್ಷರಮಾಲೆಯಲ್ಲಿ ಬ, ಸ, ವ ಎಂಬೀ ಅಕ್ಷರಗಳನ್ನು ಸುಲಭವಾಗಿ ಗುರ್ತಿಸಬಹುದು. ಹೀಗೆ ಅಕ್ಷರಗಳ ಸಮಂಜಸವಾದ ಜೋಡಣೆಯಿಂದಾಗಿರುವುದೇ “ಪದ”.\par

\indent ಮನಸ್ಸಿಗೆ ಬಂದಂತೆ ಪದಗಳನ್ನು ಕೂಡಿಸಿದರೆ ವಾಕ್ಯವಾಗುವುದಿಲ್ಲ. ಅಂತೆಯೇ ಮನಸ್ಸಿಗೆ ಬಂದಂತೆ ಅಕ್ಷರಗಳನ್ನು ಜೋಡಿಸಿದರೆ ಪದವಾಗುವುದಿಲ್ಲ.\par

\kbvexercise{ಅಭ್ಯಾಸ ಪತ್ರಿಕೆ ೧}

\vspace{2.5pt plus 3pt minus 0.5pt}\noindent\hspace*{1.2em}\kbvnum{(1)}\hspace{0.55em}ಎರಡು ಅಕ್ಷರಗಳ ಹತ್ತು ಪದಗಳನ್ನು ಬರೆಯಿರಿ.\par

\noindent\hspace*{1.2em}\kbvlabel{ಉದಾ}\hspace{0.6em}ಏತ, ಆಳ\par

\vspace{2.5pt plus 3pt minus 0.5pt}\noindent\hspace*{1.2em}\kbvnum{(2)}\hspace{0.55em}ಮೂರು ಅಕ್ಷರಗಳ ಹತ್ತು ಪದಗಳನ್ನು ಬರೆಯಿರಿ.\par

\noindent\hspace*{1.2em}\kbvlabel{ಉದಾ}\hspace{0.6em}ಅಗಸ, ಸರಳ\par
